\subsection{Первые наблюдения и ключевые эксперименты}\label{subsec:first_observations}

Долгое время считалось, что в оптических волокнах из плавленого кварца нелинейно-оптические процессы второго порядка невозможны. Кварцевое стекло является центросимметричным на макроскопическом уровне в силу своей аморфной структуры (хотя кристаллический кварц обладает тригональной сингонией). Наблюдавшуюся слабую генерацию второй гармоники в таких средах объясняли поверхностными эффектами на границе сердцевина--оболочка и квадрупольной нелинейной поляризуемостью. Помимо электродипольного механизма, существуют также магнитный дипольный, электрический квадрупольный и поверхностный вклады, однако суммарная эффективность н-о преобразования, основанная на этих эффектах, согласно расчётам не превышает $10^{-3}\%$ \cite{terhune1987second}. Таким образом, за эффективное протекание процесса отвечают какие-то другие физические механизмы.


Существенный прогресс в понимании природы н-о эффектов второго порядка в центросимметричных средах наступил в 1986 году, когда Остерберг и Маргулис сообщили об увеличении эффективности генерации второй гармоники при длительном распространении излучения накачки по стандартному одномодовому ОВ, легированному германием. \cite{osterberg1986dye}. В течение 10~часов мощность второй гармоники возрастала практически экспоненциально, а затем выходила на насыщение. Эффективность преобразования достигала $3 \%$. В качестве накачки использовалось излучение твердотельного $\text{YAG:Nd}^{3+}$ лазера с длиной волны $\lambda = 1064~\text{нм}$, модуляцией добротности и пиковой мощностью до 20~кВт. Процесс роста мощности излучения второй гармоники ($\lambda = 532~\text{нм}$) получил название самоприготовления световода.

Работа Остерберга и Маргулиса дала старт интенсивным исследованиям в данной области. В 1987 году, основываясь на экспериментах с волокнами, легированными различными примесями, авторы работы \cite{farries1987second} постулировали механизм ГВГ как самоорганизующегося процесса формирования пространственной решётки дипольных центров окраски с периодом, соответствующим эффективной генерации второй гармоники. Также впервые была измерена кривая фазового синхронизма процесса и таким образом обнаружена спектральная зависимость эффективности генерации в <<самоприготовленном>> волокне. К настоящему времени общепризнано, что эффективная генерация второй гармоники в волокне обусловлена формированием пространственной решётки квадратичной нелинейности. Однако физический механизм возникновения такой решётки остаётся предметом активных дискуссий.

Важным этапом в развитии представлений о физических механизмах н-о эффектов в волокнах являлась работа Столена и Тома \cite{stolen1987self}, согласно которой знакопеременная решётка нелинейной восприимчивости второго порядка возникала за счёт ориентации дефектов на нулевой частоте (в постоянном электрическом поле). Постоянная же поляризация, а вследствие неё и постоянное поле, появляется в результате параметрического процесса третьего порядка, в котором излучение накачки и второй гармоники (возникающей из оптического шума либо привнесённой извне). В этом случае нелинейность второго порядка будет периодичной по пространству с периодом, соответствующим периоду синхронизма для генерации второй гармоники, а величина нелинейности (и эффективность преобразования) будет пропорциональна квадрату поля волны накачки и полю волны второй гармоники:

Статическая поляризация, возникающая при смешении волн на кубической нелинейности, имеет вид:

\begin{equation}\label{eq:stolen_pdc}
	P_{dc} = \frac{3}{8} \varepsilon_0 \chi^{(3)}(0 = \omega + \omega - 2\omega) |E_\omega|^2 E_{2\omega}^* \cos(\Delta k z),
\end{equation}

где $\Delta k = k_{2\omega} - 2k_\omega$~--- фазовая расстройка. Соответствующее электростатическое поле $E_{dc} = P_{dc}/(\varepsilon_0 \chi^{(1)})$ оказывается периодическим в пространстве с периодом $\Lambda = 2\pi/\Delta k$, что обеспечивает фазовый квазисинхронизм для ГВГ. Величина коэффициента квадратичной нелинейности при этом пропорциональна квадрату поля волны накачки и полю волны второй гармоники:

\begin{equation}\label{eq:stolen_chi2}
	\chi^{(2)} \propto |E_\omega|^2 E_{2\omega}^* \cos(\Delta k z + \varphi_p).
\end{equation}

Основываясь на зависимости нелинейной восприимчивости второго порядка от поля второй гармоники \cref{eq:stolen_chi2}, авторы работы~\cite{stolen1987self} предложили и осуществили эксперимент по генерации второй гармоники, вводя в оптическое волокно затравочное излучение на удвоенной в н-о кристалле частоте. При этом время установления эффективного преобразования уменьшилось с 12~часов до~5 минут. Как оказалось впоследствии, модель Столена и Тома не может описывать процесс ГВГ, так как максимальная величина постоянного поля возникающая при н-о процессах на кубической нелинейности, по порядку величины составляет $1~$В/см и недостаточна для эффективного протекания процесса~\cite{bergot1988generation}. Тем не менее, ключевой шаг в направлении понимания физики процесса был сделан (дальнейшие модели лишь показывают способы многократного усиления поля постоянной поляризации), а успех экспериментального подтверждения теоретических предположений, сделанных авторами в начале работы, направил дальнейшие усилия многих исследователей в сторону ориентационных моделей. 

В 1988~году Уллет с соавторами показали, что записанная решётка может быть стёрта оптически с использованием излучения второй гармоники, причём затухание не было экспоненциальным \cite{ouellette1988light}. Авторы предложили зарядовую модель, согласно которой  разделение зарядов возникает в постоянном электрическом поле, образованном параметрическим смешением частот накачки и второй гармоники. В отличие от модели Столена и Тома, где предполагалась ориентация диполей, зарядовая модель Уллета объясняла нелинейность через макроскопическое разделение зарядов и возникающее электростатическое поле. Это поле, согласно соотношению $\chi^{(2)} \sim \chi^{(3)} E_{dc}$, и создавало эффективную квадратичную восприимчивость.

Таким образом, в течение первых двух лет после обнаружения эффективной ГВГ в ОВ, были проведены основные экспериментальные исследования и предложены два класса физических моделей: зарядовые и ориентационные. Обе модели корректно описывали некоторые экспериментальные факты, однако вопрос о том, какая модель реализуется на практике, оставался открытым.

